\begin{question}

  Using the orthonormality of $\ket +$ and $\ket -$, prove

  \begin{questionpart}
    $[S_i, S_j] = i \epsilon_{ijk} \hbar S_k$
  \end{questionpart}

  \begin{questionpart}
    $\{S_i, S_j\} = \left( \frac{\hbar^2}{2} \right) \delta_{ij}$
  \end{questionpart}

  \noindent where

  \begin{equation}
    \tag{1.111a}
    S_x = \frac \hbar 2 (\ket - \bra + + \ket + \bra -)
  \end{equation}

  \begin{equation}
    \tag{1.111b}
    S_y = i \frac \hbar 2 (\ket - \bra + - \ket + \bra -)
  \end{equation}
  
  \begin{equation}
    \tag{1.90}
    S_z = \frac \hbar 2 (\Lambda_+ - \Lambda_-) = (\ket + \bra + - \ket - \bra -)
  \end{equation}
\end{question}

\begin{answer}

  Strictly speaking, we'll need to show that all 9 permutations of
  $ijk$ for each part.

  %%%%%%%%%%%%%%%%%%%%%%%%%%%%%%%%%%%%%%%%%%%%%%%%%%%%%%%%%%%%%%%%%%%%%%%%%%%%%%%
  %%                              Part a                                       %%
  %%%%%%%%%%%%%%%%%%%%%%%%%%%%%%%%%%%%%%%%%%%%%%%%%%%%%%%%%%%%%%%%%%%%%%%%%%%%%%%

  \begin{answerpart}{Commutation}

    Starting with the case of $i=j$, it's pretty easy to see that

    \begin{equation}
      [S_i, S_i] = S_i^2 - S_i^2 = 0
    \end{equation}

    Next we examine the $i \ne j$ cases, of which there are 6:
    \begin{tabular}{ccc}
      $x, y$ & $y, z$ & $z, x$\\
      $y, x$ & $z, y$ & $x, z$\\
    \end{tabular}

    Since $[X, Y] = -[Y, X]$, we can trivially show that if the even (or
    cyclic) permutations are correct, then the odd (or anticyclic)
    permutations are also correct owing to the negative value
    associated with the odd (or anticyclic) terms in the Levi-Civita
    operator.

    Unfortunately, I don't see any way to, in general, show this
    result using the material available in Chapter 1.  So...let's
    crank it out.  For the 3 cyclic cases listed above, we'll need to show that:

    \begin{equation}
      [S_i, S_j] = i \hbar S_k |_{(i, j, k) \in \{(x, y, z),  (y, z, x),  (z, x, y)\}}
    \end{equation}
    
    Starting with $x$ and $y$:
    \begin{equation}
      \inlabel{a-xy}
      \begin{split}
        [S_x, S_y]
        &= S_xS_y - S_yS_x \\
        &= [\frac \hbar 2 (\ket - \bra + + \ket + \bra -)][i \frac \hbar 2 (\ket - \bra + - \ket + \bra -)] - [i \frac \hbar 2 (\ket - \bra + - \ket + \bra -)][\frac \hbar 2 (\ket - \bra + + \ket + \bra -)] \\
        &= i \frac{\hbar^2}{4} \{ [(\ket - \bra + + \ket + \bra -)][(\ket - \bra + - \ket + \bra -)] - [(\ket - \bra + - \ket + \bra -)][(\ket - \bra + + \ket + \bra -)]\}\\
        &= i \frac{\hbar^2}{4} \{ (-\Lambda_- + \Lambda_+) - (\Lambda_- - \Lambda_+) \} \\
        &= i \frac{\hbar^2}{4} \{ (-\Lambda_- + \Lambda_+) + (-\Lambda_- + \Lambda_+) \} \\
        &= i \frac{\hbar^2}{4} \{ 2 (\Lambda_+ - \Lambda_-) \} \\
        &= i \hbar \{ \frac{\hbar}{2} (\Lambda_+ - \Lambda_-) \} \\
        &= i \hbar S_z \qed{}\\
      \end{split}
    \end{equation}

    Next up to bat is $y$ and $z$:
    \begin{equation}
      \inlabel{a-yz}
      \begin{split}
        [S_y, S_z]
        &= S_yS_z - S_zS_y \\
        &= [i \frac \hbar 2 (\ket - \bra + - \ket + \bra -)][\frac \hbar 2 (\Lambda_+ - \Lambda_-)] - [\frac \hbar 2 (\Lambda_+ - \Lambda_-)][i \frac \hbar 2 (\ket - \bra + - \ket + \bra -)]\\
        &= i \frac{\hbar^2}{4} \{ [(\ket - \bra + - \ket + \bra -)][(\Lambda_+ - \Lambda_-)] - [(\Lambda_+ - \Lambda_-)][(\ket - \bra + - \ket + \bra -)] \}\\
        &= i \frac{\hbar^2}{4} \{ (\ket - \bra + + \ket + \bra -) - (-\ket - \bra + - \ket + \bra -) \}\\
        &= i \frac{\hbar^2}{4} \{ (\ket - \bra + + \ket + \bra -) + (\ket - \bra + + \ket + \bra -) \}\\
        &= i \frac{\hbar^2}{4} \{ 2(\ket - \bra + + \ket + \bra -) \}\\
        &= i \hbar \{ \frac{\hbar}{2} (\ket - \bra + + \ket + \bra -) \}\\
        &= i \hbar S_x \qed\\
      \end{split}
    \end{equation}

    Finally, we have $z$ and $x$:
    \begin{equation}
      \inlabel{a-zx}
      \begin{split}
        [S_z, S_x]
        &= S_zS_x - S_xS_z \\
        &= [\frac \hbar 2 (\Lambda_+ - \Lambda_-)][\frac \hbar 2 (\ket - \bra + + \ket + \bra -)] - [\frac \hbar 2 (\ket - \bra + + \ket + \bra -)][\frac \hbar 2 (\Lambda_+ - \Lambda_-)]\\
        &= \frac{\hbar^2}{4} \{ [(\Lambda_+ - \Lambda_-)][(\ket - \bra + + \ket + \bra -)] - [(\ket - \bra + + \ket + \bra -)][(\Lambda_+ - \Lambda_-)]\\
        &= \frac{\hbar^2}{4} \{ (-\ket - \bra + + \ket + \bra -) - (\ket - \bra + - \ket + \bra -) \}\\
        &= \frac{\hbar^2}{4} \{ 2(-\ket - \bra + + \ket + \bra -) \}\\
        &= \hbar \{ \frac{-1}{i}S_y \}\\
        &= i \hbar S_y \qed\\
      \end{split}
    \end{equation}

    In equations \inref{a-xy}, \inref{a-yz}, and \inref{a-zx} we
    showed that the cyclic scenarios are correct and the anti-cyclic
    scenarios were also shown to be correct if the cyclic scenarios
    are.
    
  \end{answerpart}


  %%%%%%%%%%%%%%%%%%%%%%%%%%%%%%%%%%%%%%%%%%%%%%%%%%%%%%%%%%%%%%%%%%%%%%%%%%%%%%%
  %%                              Part b                                       %%
  %%%%%%%%%%%%%%%%%%%%%%%%%%%%%%%%%%%%%%%%%%%%%%%%%%%%%%%%%%%%%%%%%%%%%%%%%%%%%%%

  \begin{answerpart}{Anticommutation}

    In equations \inref{a-xy}, \inref{a-yz}, and \inref{a-zx} we can
    see that each has an intermediate step in which $x - (-x)
    \Rightarrow 2x$ occurs.  Since the anticommutation would change
    that to become $x + (-x) \Rightarrow 0$ it is trivial to show that
    all scenarios in which $i \ne j$ result in 0.

    Examining the cases of $i=j$ we see that

    \begin{equation}
      \inlabel{szSquared}
      \begin{split}
        S_z^2 &= \frac{\hbar^2}{4} (\Lambda_+ - \Lambda_-)^2\\
        &= \frac{\hbar^2}{4} (\Lambda_+ + \Lambda_-)\\
        &= \frac{\hbar^2}{4} I\\
      \end{split}
    \end{equation}

    \begin{equation}
      \inlabel{sxSquared}
      \begin{split}
        S_x^2 &= \left( \frac \hbar 2 (\ket - \bra + + \ket + \bra -) \right)^2\\
        &= \frac{\hbar^2}{4}(\braket{-} + \braket{+})\\
        &= \frac{\hbar^2}{4}(\Lambda_- + \Lambda_+)\\
        &= \frac{\hbar^2}{4}\\
        &= S_z^2\\
      \end{split}
    \end{equation}

    \begin{equation}
      \inlabel{sySquared}
      \begin{split}
        S_y^2 &= \left( i \frac \hbar 2 (\ket - \bra + - \ket + \bra -) \right)^2\\
        &= - \frac{\hbar^2}{4}(-\braket - - \braket + +)\\
        &= \frac{\hbar^2}{4} (\Lambda_- + \Lambda_+) \\
        &= \frac{\hbar^2}{4} I\\
        &= S_z^2\\
      \end{split}
    \end{equation}

    Putting it all together, we see that

    \begin{equation}
      \begin{split}
        \{S_i, S_i\}
        &= S_i^2 + S_i^2\\
        &= 2S_i^2\\
        &= 2\frac{\hbar^2}{4}\\
        &= \frac{\hbar^2}{2} \qed\\
      \end{split}
    \end{equation}

    While not strictly necessary for this problem, I had these notes
    laying around so I've included them.  Add them together (as in eq
    1.116) and we have rederived equation 1.117 from the book:

    \begin{equation}
      \tag{1.116}
      S^2 \equiv S_x^2 + S_y^2 + S_z^2
    \end{equation}
    
    \begin{equation}
      \tag{1.117}
      S^2 = \left(\frac 3 4\right) \hbar^2
    \end{equation}

    Importantly, this outcome MUST be the case.  When applied with a
    test-bra/ket ($\ket \psi$), we find that:

    \begin{equation}
      \begin{split}
        \expval{S_i^{2}}{\psi}
        &= (\bra \psi S_i ) (S_i \ket \psi)\\
        &= (\alpha^* \bra{\psi'}) (\alpha \ket{\psi'})\\
        &= |\alpha^2| \braket{\psi'} \\
        &= |\alpha^2|\\
      \end{split}
    \end{equation}

    Since the magnitude of a spin-operator acting on a spin-state is
    $\frac \hbar 2$, we can quickly see that a magnitude $\frac \hbar
    2$ bra acting on its own hermitian adjoint ket must produce a
    quantity of $\frac{\hbar^2}{4}$.

  \end{answerpart}

\end{answer}
